\section{The Word Made Flesh and the Wonderful Exchange}\label{sec:theology}

A sentimental Christmas softens the Incarnation into a winter tale about a
gentle child. The 1962 Roman liturgy never allows that softening. From Midnight
Mass to the Octave Day, the Church stares straight into an abyss of paradox: the
eternal Word by whom the stars were made lies unable to speak in a feeding
trough, and the Maker of man takes mortal flesh so that mortal men may share the
life of God.

\subsection{Why the Angels Sang of Peace}

Preaching in St.~Mary Major on Christmas Day at the close of the sixth century,
Pope St.~Gregory the Great asked a question in Homily~8 on the Gospels that few
hearers ever pause to ask: why did the heavenly host cry out ``on earth peace to
men of good will'' (Luke~2:14) at Bethlehem? Gregory found the answer by
comparing the Old Testament with the Apocalypse. Before the Redeemer was born in
the flesh, fallen humanity stood in discord with the angels, estranged from
their brightness and purity by original sin and daily transgression. Because
men were strangers to God, the citizens of heaven treated them as outsiders.
When Lot fell down before two angels at Sodom (Gen~19:1) and Joshua fell on his
face before the prince of the Lord's host outside Jericho (Jos~5:15), neither
angel forbade the prostration.

Then comes the Incarnation, and everything between angels and men turns upside
down. In the Apocalypse, when St.~John falls down at the feet of the angel who
showed him the heavenly vision, the angel stops him at once: ``See thou do it
not. For I am thy fellow servant, and of thy brethren the prophets''
(Apoc~22:8--9). Gregory asks why angels accepted prostration from men before the Redeemer's
coming and shrink from it afterward: because once they behold our human
nature---which they had previously looked down upon in its weakness---lifted
above themselves in the King of heaven, they dare no longer despise beneath
their feet the humanity they venerate enthroned above them.

Because the King of heaven took to himself the clay of our flesh, the angels
recognize us once more as fellow citizens. That single sentence from Gregory
measures the height of Christmas: Bethlehem is not merely God descending to our
poverty, but humanity lifted above the seraphim in the Person of the Son.

\subsection{The Wonderful Exchange in East and West}

From the earliest Christian centuries, Greek and Latin Fathers gave this mystery
a single name: the wonderful exchange, \textit{admirabile commercium}. In the
second century, St.~Ignatius of Antioch wrote to the Ephesians (Chapter~19) that
the virginity of Mary, her childbearing, and the death of the Lord were ``three
mysteries of renown, which were wrought in silence by God,'' hidden from the
prince of this world until the new star broke the dominion of magic, ignorance,
and death. A generation later, St.~Irenaeus of Lyons taught in \textit{Against
Heresies} (Book~III, Chapter~18.1) that when the eternal Son became incarnate,
``He commenced afresh the long line of human beings, and furnished us, in a
brief, comprehensive manner, with salvation; so that what we had lost in
Adam---namely, to be according to the image and likeness of God---that we might
recover in Christ Jesus.'' In Book~III, Chapter~22.4, Irenaeus set Mary's
obedience directly beside Eve's disobedience: the knot tied by the first virgin
was untied by the Virgin Mother of the Lord.

In fourth-century Alexandria, St.~Athanasius compressed the whole economy into
ten words in \textit{On the Incarnation of the Word} (54.3): ``For He was made
man that we might be made God; and He manifested Himself by a body that we might
receive the idea of the unseen Father; and He endured the insolence of men that
we might inherit immortality.'' In Constantinople on Christmas Day~380,
St.~Gregory of Nazianzus broke into wonder in Oration~38 (section~13):

\begin{quote}
O new commingling; O strange conjunction; the Self-Existent comes into being,
the Uncreate is created, That which cannot be contained is contained, by the
intervention of an intellectual soul, mediating between the Deity and the
corporeity of the flesh. And He Who gives riches becomes poor, for He assumes
the poverty of my flesh, that I may assume the richness of His Godhead. He that
is full empties Himself, for He empties Himself of His glory for a short while,
that I may have a share in His Fulness.
\end{quote}

In the Latin West, St.~Ambrose of Milan read the details of Luke~2:7 as one
unbroken barter of mercy in his \textit{Exposition of the Gospel according to
Luke} (Book~II, section~41). He was an infant, Ambrose writes, so that you might
grow into a full-grown man; wrapped in swaddling bands, so that you might be
loosed from the snares of death; laid in a manger, so that you might be fed at
the altar (\textit{ille in praesepibus, ut tu in altaribus}); without room in the inn, so that you might have many mansions in heaven.
Therefore, Ambrose concludes, Christ's poverty becomes our inheritance, and the
weakness of the Lord becomes our strength. Three centuries later in Northumbria, St.~Bede the
Venerable echoed Ambrose line for line in his commentary on Luke~2:7: ``He Who
is the bread of Angels is laid down in a manger, that He might feast us, as it
were the sacred animals, with the bread of His flesh.''

St.~Augustine brought that exchange down to a wake-up call for drowsy hearts.
In Sermon~185 (section~1), he cries out: \textit{Expergiscere, homo: pro te Deus
factus est homo}, ``Wake up, O man: for you God was made man.'' Truth, which
abides in the bosom of the Father, sprang from the earth to lie in a mother's
lap; Truth, which the heavens cannot hold, was placed in a manger. You would
have stayed dead forever had he not been born in time; \textit{Perisses, nisi
venisset}, ``You would have perished, had he not come.'' In Sermon~191
(section~1), Augustine piles paradox upon paradox: \textit{Homo factus, hominis
factor}, ``Made man, the Maker of man,'' so that he who rules the stars might
nurse at a mother's breast, Bread might hunger, the Fountain might thirst, Light
might sleep, the Way might grow weary on the road, and Life itself might die. In
his second Tractate on John (section~16), commenting on ``the Word was made
flesh'' (John~1:14), Augustine explains why God healed us from within our own
flesh:

\begin{quote}
By dust thou wert blinded, and by dust thou art healed: flesh, then, had wounded
thee, flesh heals thee. The soul had become carnal by consenting to the
affections of the flesh; thus had the eye of the heart been blinded. ``The Word
was made flesh:'' that Physician made for thee an eye-salve.
\end{quote}

In fifth-century Ravenna, St.~Peter Chrysologus drew the moral consequence in
Sermon~148: that God is found in our flesh is no injury to the Creator, but the
supreme honor of the creature. Why, he asks his hearers, should man think
himself cheap in his own eyes when he is so precious to God, or dishonor himself
when God has so honored him?

\subsection{True God and True Man: Pope St.~Leo the Great}

No preacher shaped the Roman Christmas liturgy more deeply than Pope St.~Leo the
Great, whose Nativity sermons are read year after year in the Roman Breviary.
Facing both Manichaean contempt for matter and Eutychian confusion of Christ's
two natures, Leo insisted in Sermon~22 (sections~1--2) that true God was born in
the whole and perfect nature of true man---complete in what is his own, complete
in what is ours, taking our nature without our sin. His self-emptying was ``the
condescension of Pity not the failing of Power,'' entering our lower world ``in
a new order, by a new nativity,'' invisible in his own nature yet visible in
ours, abiding before all time yet beginning to be in time. In Sermon~23
(sections~2--3), Leo binds the Nativity directly to the Cross: both natures keep
their own character without loss, because a Victim had to be offered for our
atonement who was at once a true partner of our race and wholly free from our
stain.

Because the Nativity and the Passion form one saving design, Leo closes his
First Sermon on the Nativity (Sermon~21, section~3---Lesson~6 of Christmas
Matins) with the sentence that echoes through every hour of Christmastide:

\begin{quote}
Christian, acknowledge thy dignity, and becoming a partner in the Divine nature,
refuse to return to the old baseness by degenerate conduct. Remember the Head
and the Body of which thou art a member. Recollect that thou wert rescued from
the power of darkness and brought out into God's light and kingdom.
\end{quote}

\subsection{The Mother of God, the Manger, and the Altar}

The Church never separates the reality of the Incarnation from the Virgin who
bore the Child or from the Eucharist that prolongs his presence. In
eighth-century Palestine, St.~John of Damascus wrote in \textit{An Exact
Exposition of the Orthodox Faith} (Book~III, Chapter~12) that the title
\textit{Theotokos}, Mother of God, ``embraces the whole mystery of the
dispensation,'' because it confesses in a single word the one Person, the two
natures, and the two births---eternal and temporal---of our Lord Jesus Christ.
In Book~IV, Chapter~11, he numbers ``His sacred tabernacles which are the
manger, the cave, Golgotha'' among the places venerated by Christians because
God truly entered material space and sanctified it.

In fifth-century Alexandria, St.~Cyril looked at the manger in his first Sermon
on Luke and saw the altar of the Church: finding man reduced to the level of
beasts, Christ placed himself like fodder in a feeding trough so that, leaving
behind a brutish life, ``by now approaching the manger, even His own table, we
find no longer fodder, but the bread from heaven, which is the body of life.''
St.~Maximus of Turin (preserved in St.~Thomas Aquinas's \textit{Catena Aurea} on
Luke~2:12) answered anyone tempted to stumble over such poverty: ``If thou
despisest the manger, raise thy eyes a little, and behold the new star in heaven
proclaiming to the world the Lord's nativity.''

Medieval doctors gathered those patristic threads into serene clarity.
Preaching in his First Sermon for Christmas Day, St.~Bernard of Clairvaux
marveled that the Father had ``shortened His own Word'': he who fills heaven and
earth lies in a narrow crib, and he who is from eternity becomes an infant of a
day. ``It is intolerable impudence on the part of a wretched little worm to
inflate and enlarge itself, when the Lord of glory has voluntarily laid aside
His greatness.'' In his Fifth Christmas Sermon, Bernard added that ``the infant
Word of God, Whose very infancy is divinely eloquent,'' speaks through the
stable, the crib, the tears, and the swaddling bands to shelter and bind the
wounds of the man who fell among robbers.

In the \textit{Summa Theologiae} (Tertia Pars, q.~1, a.~1), St.~Thomas Aquinas
reduced the whole mystery to the nature of goodness itself: ``it belongs to the
essence of the highest good to communicate itself in the highest manner to the
creature,'' which God accomplished by uniting created nature to himself in one
Person. Why Bethlehem rather than imperial Rome or royal Jerusalem? In Question
35, Article~7, Thomas answers with the Council of Ephesus: had Christ chosen
great Rome, men would have credited the transformation of the world to human
power; instead he chose a poor mother and a small town, willing ``to be born in
a mean city, and to suffer reproach in a great city,'' so that the renewal of
the earth might stand unmistakably as the work of God.
